\section{Conclusion}

We presented \eva{}, a production-grounded testing design for assessing agentic models under realistic infrastructure engineering conditions.
Unlike traditional benchmarks that provide fixed datasets and scoring functions, \eva{} provides a \emph{methodology}---a unified assessment architecture supporting heterogeneous LLM providers and agent SDKs, realistic long-horizon compiler construction workloads with strict serial dependencies, a \resurrection{} mechanism that decomposes cascading failure into diagnostic components, and utility-oriented multi-dimensional metrics measuring both outcome quality and process efficiency.

By leveraging authentic workloads from the \yatcc{} platform, our framework continuously assesses frontier models through genuine development tasks spanning lexical analysis, syntax parsing, IR generation, optimization, and backend code generation.
Our experiments across 22 model configurations, four agent backends, and two workload difficulty levels demonstrate that:
(1)~serial evaluation preserves overall rankings ($\tau = 0.763$) but reveals specific failure patterns invisible to static benchmarks;
(2)~resurrection recovers substantial downstream diagnostic signal (median 73.65-point pipeline score recovery, 17/19 configurations with $>40$-point improvement);
(3)~process cost is decoupled from outcome ($R^2 = 0.007$), providing diagnostic rather than predictive value;
(4)~the difficulty gap between \yatcc{} and \yatcc{}-Hard (up to 34.4-point decline for top models) quantifies the value of scaffolding and reveals which models depend on provided structure versus which can operate autonomously.

These findings suggest that scaffold design and evaluation methodology may become as important as foundation model capability itself in determining real-world agent performance.
\eva{} thus establishes a new paradigm for agent evaluation that moves beyond synthetic, isolated benchmarks toward continuous, runtime-observable, production-grounded assessment.
We release \eva{} as an open-source framework with complete evaluation infrastructure, scoring scripts, containerized execution environments, and reproducible workflows.